% Generated draft: requires mathematical and source review.
\documentclass[11pt]{article}
\usepackage[margin=25mm]{geometry}
\usepackage{fontspec}
\setmainfont{Latin Modern Roman}
\usepackage{amsmath,amssymb,mathtools,unicode-math}
\setmathfont{Latin Modern Math}
\usepackage{enumitem,needspace,xurl,hyperref}
\setlength{\parindent}{0pt}
\setlength{\parskip}{5pt}
\setlength{\emergencystretch}{3em}
\newcommand{\N}{\mathbb{N}}
\newcommand{\Z}{\mathbb{Z}}
\newcommand{\Q}{\mathbb{Q}}
\newcommand{\R}{\mathbb{R}}
\newcommand{\C}{\mathbb{C}}
\newcommand{\F}{\mathbb{F}}
\newcommand{\A}{\mathbb{A}}
\newcommand{\PP}{\mathbb{P}}
\newcommand{\E}{\mathbb{E}}
\newcommand{\eps}{\varepsilon}
\newcommand{\dd}{\,\mathrm{d}}
\begin{document}
\section*{UnsolvedMath ID: 30003239}
\textbf{Source problem code:} catalogue entry
\subsection{Definitions and mathematical statement}
Let \(k\) be a number field, let \(\Omega_k\) be its set of places, and write \(\A_k\) for its ring of adeles. For a proper variety \(Y\) over a number field \(L\), put
\[
\operatorname{Br}(Y)=H^2_{\mathrm{\acute et}}(Y,\mathbb G_m)_{\mathrm{tors}},
\qquad
Y(\A_L)^{\operatorname{Br}(Y)}
=
\left\{(y_v)\in Y(\A_L):
\sum_{v\in\Omega_L}\operatorname{inv}_v\bigl(\alpha(y_v)\bigr)=0
\text{ for every }\alpha\in\operatorname{Br}(Y)\right\}.
\]
Here \(\alpha(y_v)\in\operatorname{Br}(L_v)\) is obtained by evaluation, and \(\operatorname{inv}_v\) is the local invariant map. Saying that the Brauer--Manin obstruction is the only obstruction to the Hasse principle for \(Y/L\) means
\[
Y(\A_L)^{\operatorname{Br}(Y)}\ne\varnothing
\quad\Longrightarrow\quad
Y(L)\ne\varnothing.
\tag{H}
\]
This is an existence assertion; it does not assert density or weak approximation.

Let
\[
f:X\longrightarrow\mathbb P^1_k
\]
be a projective surjective morphism whose generic fibre \(X_\eta\), over \(K=k(\mathbb P^1_k)\), is geometrically rationally connected. This last condition means rational connectedness after extension to an algebraic closure \(\overline K\); in the usual smooth projective setting, two general geometric points can be joined by the image of a morphism from \(\mathbb P^1_{\overline K}\).

\textbf{Question supplied by the record.} If every smooth fibre of \(f\) satisfies \((\mathrm H)\), must \(X/k\) satisfy \((\mathrm H)\)?

For the interpretation involving fibres over \(k\)-rational base points, define
\[
T_k=\{t\in\mathbb P^1(k):X_t=f^{-1}(t)\text{ is smooth over }k\}.
\]
The question then has the explicit form
\[
\left[
\forall t\in T_k,\quad
X_t(\A_k)^{\operatorname{Br}(X_t)}\ne\varnothing
\Longrightarrow X_t(k)\ne\varnothing
\right]
\quad\Longrightarrow\quad
\left[
X(\A_k)^{\operatorname{Br}(X)}\ne\varnothing
\Longrightarrow X(k)\ne\varnothing
\right].
\]
\textbf{Unresolved qualifications.} The source does not specify whether ``every smooth fibre'' means only fibres over \(k\)-rational points or all smooth fibres over closed points. In the latter interpretation, the hypothesis is instead \((\mathrm H)\) for \(X_t/k(t)\) at every closed point \(t\) with smooth fibre, where \(k(t)\) is its residue number field. These quantifiers are different, and the intended one requires review. The source also does not explicitly assume that \(X\) is smooth or geometrically integral, or specify its scheme-theoretic conventions. No such additional hypothesis is inserted here; the intended total-space hypotheses and the convention for rational connectedness must be confirmed before this becomes a fully qualified formulation.
\subsection{Short English statement}
The supplied question asks whether the sufficiency of the Brauer--Manin obstruction for rational points on every smooth fibre passes to the total space of a projective fibration over \(\mathbb P^1\) with geometrically rationally connected generic fibre. The scope of the fibre quantifier and the regularity and integrality hypotheses on the total space remain unresolved. The supplied status is partially solved; this is a statement-only draft requiring human mathematical review.
\subsection{Sources}
\begin{itemize}
\item Supplied local problem registry record: permanent UnsolvedMath ID \(30003239\), exact source problem code \texttt{catalogue entry}, title \emph{Brauer--Manin Obstructions for Rationally Connected Fibrations}, in the Oberwolfach Reports Open Problems collection. No individual proposer is supplied; the proposed year is recorded as 2016.
\item Supplied report citation: \emph{Definability and Decidability Problems in Number Theory} (2017), Oberwolfach Reports, pp.~2793--2866, DOI \texttt{10.4171/owr/2016/49}. \url{https://doi.org/10.4171/owr/2016/49}. The record attributes the question to workshop open-problem discussions collected in this report; no problem-specific page or named contributor is supplied.
\item Harpaz--Wittenberg, ``fibration method for zero-cycles and rational points'' (supplied reference label). \url{https://doi.org/10.4007/annals.2016.183.1.5}.
\item Browning--Matthiesen, ``norm forms'' (supplied reference label). \url{https://doi.org/10.24033/asens.2348}.
\item Harpaz--Wei--Wittenberg, ``few non-split fibres'' (supplied reference label). \url{https://doi.org/10.1515/crelle-2022-0042}.
\item Literature assessment dated exactly \texttt{2026-08-21}: supplied dated curation, not independently verified. It labels the problem partially solved and states that Harpaz--Wittenberg prove the zero-cycle analogue for rationally connected fibrations and rational-point results under a polynomial-value conjecture; unconditional rational-point density is known in substantial norm-form and few-fibre cases, while the full question remains open. The background identifies this as literature triage derived from the supplied OWR curation and requests independent verification of time-sensitive status and cited claims.
\item The registry supplies the statement label \texttt{corrected\_verified} and reports a manual semantic, typographic, and LaTeX review claiming self-containedness and preservation of the mathematical target. These are supplied editorial assertions, not independent verification of the formulation or its mathematical status.
\end{itemize}
\end{document}
